\documentclass[10pt,a4paper]{amsart}
\usepackage[margin=19mm]{geometry}
\usepackage{amsmath,amssymb,mathtools}
\usepackage[scr=rsfs,cal=euler]{mathalfa}
\usepackage{xfrac}
\usepackage[unicode,hidelinks,bookmarks=false]{hyperref}
\newcommand{\ie}{i.e.\ }
\newcommand{\cf}{cf.\ }
\newcommand{\resp}{respectively\ }
\newcommand{\Subsection}{\subsection}
\providecommand{\nobreakdash}{\nobreak\hbox{-}}
\setlength{\emergencystretch}{2em}
\multlinegap0pt
\usepackage[shortlabels]{enumitem}
\usepackage{amssymb}
\usepackage{mathtools}
\newtheorem{thm}{}
\newtheorem{defn}{}
\newcommand{\Gcinf}{\mathcal{G}\cinfty}
\newcommand{\N}{\field{N}} % naturals
\newcommand{\R}{\field{R}} % reals
\newcommand{\RC}[1]{\frontRise{\R}{#1}\Rtil}
\newcommand{\Rtil}{\widetilde \R} % real Colombeau generalized number
\newcommand{\cinfty}{\mbox{\ensuremath{\mathcal{C}}}^{\infty}}
\newcommand{\diff}[1]{\ifmmode\mathchoice{\hbox{\rm d}#1}  % displaystyle
 {\hbox{\rm d}#1}  % normal 
 {\scalebox{0.75}{$\hbox{\rm d}#1$}}  % scriptstyle 
 {\scalebox{0.35}{$\hbox{\rm d}#1$}}  % scriptscriptstyle
 \fi} % dt,dx,... for integrals
\newcommand{\eps}{\varepsilon} % for the sake of brevity only
\newcommand{\field}[1]{\mathbb{#1}}
\newcommand{\frontRise}[2]{\ifmmode\mathchoice{{\vphantom{#1}}^{\scalebox{0.6}{$#2$}}}  % displaystyle
 {{\vphantom{#1}}^{\scalebox{0.56}{$#2$}}}  % normal 
 {{\vphantom{#1}}^{\scalebox{0.47}{$#2$}}}  % scriptstyle 
 {{\vphantom{#1}}^{\scalebox{0.35}{$#2$}}}\fi} % scriptscriptstyle 
\newcommand{\gsf}{\frontRise{\mathcal{G}}{\rho}\mathcal{GC}^{\infty}}
\newcommand{\rcrho}{\RC{\rho}}
\newcommand{\rti}{\RC{\rho}}
\title{Classical finite dimensional fixed point methods for generalized functions}
\author{Kevin Islami, George Apaaboah, Paolo Giordano}
\date{Source: arXiv:2603.08473v1 (CC BY 4.0)}
\begin{document}
\maketitle

\noindent\emph{Source and licence.} Classical finite dimensional fixed point methods for generalized functions. Version arXiv:2603.08473v1 (CC BY 4.0), distributed by arXiv under the Creative Commons Attribution 4.0 International licence, http://creativecommons.org/licenses/by/4.0/, as recorded in the arXiv licence metadata for this exact version and shown on the abstract page https://arxiv.org/abs/2603.08473v1. Full licence notice: \texttt{licenses/arxiv-2603.08473-CC-BY-4.0.txt}.\par\smallskip

\noindent\textit{Editorial scope.} Counted unit: the article's thm thm (source0.tex lines 972--979). The article's complete original proof of it is delivered with it (source0.tex lines 981--1020, one \texttt{proof} environment of 40 lines). No mathematics is written, repaired or re-derived: the statement and the proof are the article's own lines, in the article's own notation, and no retained line is substituted or re-flowed.  Retained uncounted as the source-local context the proof needs: lines 490--505; lines 509--528; lines 566--582; lines 805--814.  Nothing was omitted from any retained span, so the package carries no omission record.  No retained line contains a citation, so the wrapper carries no bibliography and no bibliography entry.  No retained line contains a figure environment, an image-inclusion command or drawing code, so no image is delivered and the package declares no asset dependency.  Editorial formatting only, in addition: an AMS \texttt{amsart} preamble that restates the article's own declarations and macros used by the retained lines, namely the environments \texttt{thm}, \texttt{defn} on separate counters, together with the standard packages those lines need.  The retained environments print in this document as Theorem 1, Definition 1 in order of appearance, whereas the article prints the same items with the numbering of its own full document.  The complete native source of the article as distributed in the arXiv e-print for this exact version is bundled unchanged as \texttt{sources/196006/source0.tex}.
\begin{thm}
\label{thm:propGSF}Let $X\subseteq\RC{\rho}^{n}$ and $Y\subseteq\RC{\rho}^{d}$
be arbitrary subsets of generalized points. Let $f_{\eps}\in\cinfty(\R^{n},\R^{d})$
be a net of smooth functions that defines a generalized smooth map
of the type $X\longrightarrow Y$, then
\begin{enumerate}
\item $\forall\alpha\in\N^{n}\,\forall(x_{\eps}),(x'_{\eps})\in\R_{\rho}^{n}:\ [x_{\eps}]=[x'_{\eps}]\in X\ \Rightarrow\ (\partial^{\alpha}f_{\eps}(x_{\eps}))\sim_{\rho}(\partial^{\alpha}f_{\eps}(x'_{\eps}))$.
\item \label{enu:GSF-cont}Each $f\in\gsf(X,Y)$ is continuous with respect
to the sharp topologies induced on $X$, $Y$.
\item $f:X\longrightarrow Y$ is a GSF if and only if there exists a net
$v_{\eps}\in\cinfty(\Omega_{\eps},\R^{d})$ such that $\left\langle \Omega_{\eps}\right\rangle :=\left\{ x\in\rti^{n}\mid\forall[x_{\eps}]=x\,\forall^{0}\eps:\ x_{\eps}\in\Omega_{\eps}\right\} \supseteq X$,
$(\partial^{\alpha}v_{\eps}(x_{\eps}))\in\R_{{\scriptscriptstyle \rho}}^{d}$
for all representatives $[x_{\eps}]=x\in X$ and all $\alpha\in\N^{n}$,
and $f=[v_{\eps}(-)]|_{X}$.
\end{enumerate}
\end{thm}

\begin{thm}[Fermat-Reyes theorem for GSF]
\label{thm:FR-forGSF} Let $U\subseteq\RC{\rho}^{n}$ be a sharply
open set, let $v=[v_{\eps}]\in\RC{\rho}^{n}$, and let $f\in\gsf(U,\RC{\rho})$
be a generalized smooth map generated by the net of smooth functions
$f_{\eps}\in\cinfty(\Omega_{\eps},\R)$. Then
\begin{enumerate}
\item \label{enu:existenceRatio}There exists a sharp neighborhood $T$
of $U\times\{0\}$ and a generalized smooth map $r\in\gsf(T,\RC{\rho})$,
called the \emph{generalized incremental ratio} of $f$ \emph{along}
$v$, such that 
\[
\forall(x,h)\in T:\ f(x+hv)=f(x)+h\cdot r(x,h).
\]
\item \label{enu:uniquenessRatio}Any two generalized incremental ratios
coincide on a sharp neighborhood of $U\times\{0\}$.
\item \label{enu:defDer}We have $r(x,0)=\left[\frac{\partial f_{\eps}}{\partial v_{\eps}}(x_{\eps})\right]$
for every $x\in U$ and we can thus define $\diff f(x)\cdot v:=\frac{\partial f}{\partial v}(x):=r(x,0)$,
so that $\frac{\partial f}{\partial v}\in\gsf(U,\RC{\rho})$.
\end{enumerate}
\end{thm}

\begin{thm}
\label{thm:Taylor}Let $f\in\gsf(U,\rcrho)$ be a generalized smooth
function defined in the sharply open set $U\subseteq\rcrho^{d}$.
Let $a$, $b\in\rcrho^{d}$ such that the line segment $[a,b]\subseteq U$,
and set $h:=b-a$. Then, for all $n\in\N$ we have
\begin{enumerate}
\item \label{enu:LagrangeRem}$\exists\xi\in[a,b]:\ f(a+h)=\sum_{j=0}^{n}\frac{\diff{^{j}f}(a)}{j!}\cdot h^{j}+\frac{\diff{^{n+1}f}(\xi)}{(n+1)!}\cdot h^{n+1}.$
\item \label{enu:integralRem}$f(a+h)=\sum_{j=0}^{n}\frac{\diff{^{j}f}(a)}{j!}\cdot h^{j}+\frac{1}{n!}\cdot\int_{0}^{1}(1-t)^{n}\,\diff{^{n+1}f}(a+th)\cdot h^{n+1}\,\diff{t}.$
\end{enumerate}
\noindent Moreover, there exists some $R\in\rcrho_{>0}$ such that
\begin{equation}
\forall k\in B_{R}(0)\,\exists\xi\in[a,a+k]:\ f(a+k)=\sum_{j=0}^{n}\frac{\diff{^{j}f}(a)}{j!}\cdot k^{j}+\frac{\diff{^{n+1}f}(\xi)}{(n+1)!}\cdot k^{n+1}\label{eq:LagrangeInfRest-1-1}
\end{equation}
\begin{equation}
\frac{\diff{^{n+1}f}(\xi)}{(n+1)!}\cdot k^{n+1}=\frac{1}{n!}\cdot\int_{0}^{1}(1-t)^{n}\,\diff{^{n+1}f}(a+tk)\cdot k^{n+1}\,\diff{t}\approx0.\label{eq:integralInfRest-1-1}
\end{equation}
\end{thm}

\begin{defn}
\label{def:NRP}Let $U\subseteq\rcrho^{d}$ and $f\in\Gcinf(U,\rcrho^{d})$.
We say that $(x_{n})_{n\in\N}$ is \emph{a Newton(-Raphson) process
for $f$} if for all $n\in\N$ we have:
\begin{enumerate}
\item $x_{n}\in U$;
\item $\diff f(x_{n})$ is invertible;
\item $x_{n+1}=x_{n}-[\diff f(x_{n})]^{-1}f(x_{n})$.
\end{enumerate}
\end{defn}

\begin{thm}
Let $U\subseteq\rti^{d}$ and $f\in\gsf(U,\rti^{d})$ be such that
at $x^{*}\in U$ the differential $\diff f(x^{*})$ is invertible.
Assume that $(x_{n})_{n\in\N}$ is a Newton process for $f$ which
converges to $x^{*}$ as $n\to\infty$. Then for all $M>\frac{1}{2}\left|\diff f(x^{*})^{-1}\diff f(x^{*})\right|$
and for $n\mathbb{\in N}$ sufficiently large, we have $|x_{n+1}-x^{*}|\le M|x_{n}-x^{*}|^{2}$,
i.e.~$(x_{n})_{n\in\N}$ converges quadratically to $x^{*}$.
\end{thm}

\begin{proof}
Set $x_{n}-x^{*}=:h_{n}$. Using Taylor's theorem \ref{thm:Taylor},
for some $\mu_{n}\in[x_{n},x^{*}]$ we have

\[
f(x_{n}-h_{n})=f(x_{n})-h_{n}\cdot\diff f(x_{n})+\frac{h_{n}^{2}}{2}\diff f(\mu_{n}).
\]

\noindent Now, since $x_{n}-h_{n}=x^{*}$ and the Newton process $(x_{n})_{n\in\N}$
converges to $x^{*}$, we necessarily have $f(x^{*})=0$.

\[
f(x_{n})-(x_{n}-x^{*})\diff f(x_{n})+\frac{h_{n}^{2}}{2}\diff f(\mu_{n})=0.
\]

\noindent By Def.~\ref{def:NRP}, $\diff f(x_{n})$ is invertible,
hence

\[
\diff f(x_{n})^{-1}f(x_{n})-(x_{n}-x^{*})+\frac{h_{n}^{2}}{2}\diff f(x_{n})^{-1}\diff f(\mu_{n})=0
\]

\noindent From Def.~\ref{def:NRP} of Newton's process, we have

\[
x_{n+1}-x^{*}=\frac{h_{n}^{2}}{2}\diff f(x_{n})^{-1}\diff f(\mu_{n})
\]

\noindent which implies

\[
|x_{n+1}-x^{*}|=\frac{1}{2}\left|\diff f(x_{n})^{-1}\diff f(\mu_{n})\right||x_{n}-x^{*}|^{2}
\]

\noindent By sharp continuity Thm.~\ref{thm:propGSF} and the Fermat-Reyes
Thm.~\ref{thm:FR-forGSF}, $\diff f(x_{n})$ converges to $\diff f(x^{*})$
and $\mu_{n}\in[x_{n},x^{*}]$ converges to $x^{*}$, and for sufficiently
large $n$, $M>\frac{1}{2}\left|\diff f(x_{n})^{-1}\diff f(\mu_{n})\right|\to\frac{1}{2}\left|\diff f(x^{*})^{-1}\diff f(x^{*})\right|$
and this proves the claim.
\end{proof}

\end{document}
